\documentclass[12pt,a4paper]{article}

% ------------------------------------------------
% Pacotes
% ------------------------------------------------
\usepackage[utf8]{inputenc}
\usepackage[T1]{fontenc}
\usepackage[brazil]{babel}
\usepackage{geometry}
\usepackage{enumitem}
\usepackage{amsmath}
\usepackage[table]{xcolor}
\usepackage{listings}
\usepackage{multicol}
\usepackage{courier}

\lstset{
    language=Java,
    keywordstyle=\bfseries,
    showstringspaces=false,
    breaklines=true,
    frame=none,
    numbers=none,
    tabsize=2,
    commentstyle=\itshape\color{gray},
    morekeywords={this, class, private, public, String, val, var, fun, data, override, return, listOf, mutableListOf, println, it, filter, map, forEach, sortedBy, firstOrNull, Modifier, Column, Row, Text, Image, Card, LazyColumn, Scaffold, Box, Spacer, Button, onClick, padding, fillMaxSize, fillMaxWidth, height, width, dp, sp},
    extendedchars=true,
    inputencoding=utf8
}

\geometry{
    a4paper,
    left=2cm,
    right=2cm,
    top=2cm,
    bottom=2cm
}

\setlength{\parindent}{0pt}

% Caixa para resposta discursiva
\newcommand{\resposta}[1]{
\noindent
\framebox[\textwidth]{
\begin{minipage}{0.98\textwidth}
\vspace*{#1}\null
\end{minipage}
}
}

\begin{document}

\begin{center}
    {\LARGE \textbf{LISTA DE EXERCÍCIOS DE REVISÃO — N1}} \\[0.3cm]
    {\large Desenvolvimento de Software para Dispositivos Móveis} \\[0.2cm]
    {\normalsize Engenharia de Software — UniRV}
\end{center}

\vspace{0.5cm}
\hrule
\vspace{0.5cm}

\textbf{Código fonte de apoio para as questões 1 até 4}

\setlength{\columnseprule}{0.4pt}
\begin{multicols}{2}
\begin{lstlisting}[basicstyle=\scriptsize\ttfamily, breaklines=true]
package com.example.livraria

class MainActivity : ComponentActivity() {
    override fun onCreate(
        savedInstanceState: Bundle?
    ) {
        super.onCreate(savedInstanceState)
        enableEdgeToEdge()
        setContent {
            LivrariaTheme {
                Scaffold(
                    modifier = Modifier
                        .fillMaxSize()
                ) { innerPadding ->
                    TelaInicial(
                        modifier = Modifier
                            .padding(innerPadding)
                    )
                }
            }
        }
    }
}

data class Livro(
    val titulo: String,
    val autor: String,
    val preco: Double,
    val avaliacao: Int,
    val capa: String
)

@Composable
fun LivroCard(
    titulo: String,
    autor: String,
    preco: Double,
    avaliacao: Int,
    capa: String,
    modifier: Modifier = Modifier
) {
    Card(
        modifier = modifier
            .fillMaxWidth()
            .padding(8.dp),
        shape = RoundedCornerShape(12.dp),
        colors = CardDefaults.cardColors(
            containerColor =
                Color(0xFFF5F5DC)
        ),
        elevation =
            CardDefaults.cardElevation(
                defaultElevation = 4.dp
            )
    ) {
        Column(
            modifier = Modifier
                .fillMaxWidth()
                .padding(12.dp),
            horizontalAlignment =
                Alignment.CenterHorizontally
        ) {
            Image(
                painter =
                    painterResource(id = capa),
                contentDescription =
                    "Capa de $titulo",
                modifier = Modifier
                    .height(180.dp)
                    .fillMaxWidth(),
                contentScale =
                    ContentScale.Fit
            )

            Spacer(
                modifier =
                    Modifier.height(8.dp)
            )

            Text(
                text = titulo,
                fontSize = 14.sp,
                fontWeight = FontWeight.Bold,
                color = Color(0xFF333333),
                textAlign = TextAlign.Center
            )

            Text(
                text = autor,
                fontSize = 12.sp,
                color = Color(0xFF777777)
            )

            Text(
                text = "R$ ${"%.2f"
                    .format(preco)}",
                fontSize = 14.sp,
                fontWeight = FontWeight.Bold,
                color = Color(0xFF2E7D32),
                modifier = Modifier
                    .padding(top = 6.dp)
            )

            Row(
                modifier = Modifier
                    .padding(top = 4.dp)
            ) {
                repeat(avaliacao) {
                    Text(
                        text = "estrela",
                        fontSize = 12.sp,
                        color =
                            Color(0xFFFFC107)
                    )
                }
            }
        }
    }
}

@Composable
fun ListaLivros() {
    val livros = listOf(
        Livro(
            titulo = "O Senhor dos Aneis",
            autor = "J.R.R. Tolkien",
            preco = 59.90,
            avaliacao = 5,
            capa = "capa_senhor"
        ),
        Livro(
            titulo = "Dom Casmurro",
            autor = "Machado de Assis",
            preco = 29.90,
            avaliacao = 4,
            capa = "capa_dom"
        ),
        Livro(
            titulo = "1984",
            autor = "George Orwell",
            preco = 39.90,
            avaliacao = 5,
            capa = "capa_1984"
        )
    )

    LazyColumn(
        modifier = Modifier.fillMaxSize(),
        contentPadding =
            PaddingValues(16.dp),
        verticalArrangement =
            Arrangement.spacedBy(12.dp)
    ) {
        items(livros) { livro ->
            LivroCard(
                titulo = livro.titulo,
                autor = livro.autor,
                preco = livro.preco,
                avaliacao = livro.avaliacao,
                capa = livro.capa
            )
        }
    }
}

@Composable
fun TelaInicial(modifier: Modifier) {
    Column(modifier = modifier) {
        Row(
            modifier = Modifier
                .fillMaxWidth()
                .padding(16.dp),
            horizontalArrangement =
                Arrangement.Center
        ) {
            Text(
                text = "Livraria ",
                fontSize = 20.sp,
                fontWeight = FontWeight.Bold
            )
            Text(
                text = "Virtual",
                fontSize = 20.sp,
                color = Color(0xFF6A1B9A)
            )
        }
        ListaLivros()
    }
}

@Preview(showBackground = true)
@Composable
fun LivrariaPreview() {
    LivrariaTheme {
        TelaInicial(
            modifier = Modifier
        )
    }
}
\end{lstlisting}
\end{multicols}

\vspace{0.5cm}

% ================================================================
% QUESTÃO 1 — @Composable
% ================================================================
\textbf{Questão 1 —} No código de apoio, diversas funções como \texttt{LivroCard}, \texttt{ListaLivros} e \texttt{TelaInicial} são marcadas com a anotação \texttt{@Composable}. Explique qual é a finalidade dessa anotação no Jetpack Compose e por que uma função comum (sem \texttt{@Composable}) não poderia ser utilizada para descrever elementos de interface.

\vspace{0.3cm}
\resposta{6cm}

\vspace{0.5cm}

% ================================================================
% QUESTÃO 2 — Row, Column, layouts
% ================================================================
\textbf{Questão 2 —} Analisando a função \texttt{TelaInicial} no código de apoio, percebe-se o uso combinado de \texttt{Column} e \texttt{Row}. Responda:

\begin{enumerate}[label=\alph*)]
    \item Qual é a diferença de comportamento entre \texttt{Column} e \texttt{Row} na organização dos elementos filhos?
    \item No trecho do \texttt{Row} que exibe o título, qual é a função do parâmetro \texttt{horizontalArrangement = Arrangement.Center}?
    \item Explique o que faz o componente \texttt{Spacer} utilizado dentro do \texttt{LivroCard} e por que ele é útil na construção de layouts.
\end{enumerate}

\vspace{0.3cm}
\resposta{7cm}

\vspace{0.5cm}

% ================================================================
% QUESTÃO 3 — dp e sp
% ================================================================
\textbf{Questão 3 —} No código de apoio, observa-se o uso frequente das unidades \texttt{dp} (ex: \texttt{8.dp}, \texttt{12.dp}, \texttt{180.dp}) e \texttt{sp} (ex: \texttt{14.sp}, \texttt{12.sp}, \texttt{20.sp}). Responda:

\begin{enumerate}[label=\alph*)]
    \item Explique a diferença entre \texttt{dp} e \texttt{sp} e indique para qual tipo de elemento cada unidade é mais adequada.
    \item Imagine que um desenvolvedor utilizou \texttt{180px} em vez de \texttt{180.dp} para a altura da imagem da capa do livro. Descreva o que aconteceria ao executar o aplicativo em dois dispositivos com densidades de tela diferentes (ex: mdpi e xxxhdpi).
\end{enumerate}

\vspace{0.3cm}
\resposta{7cm}

\vspace{0.5cm}

% ================================================================
% QUESTÃO 4 — @Preview vs. Activity
% ================================================================
\textbf{Questão 4 —} No código de apoio, existe a função \texttt{LivrariaPreview()} anotada com \texttt{@Preview(showBackground = true)} e a classe \texttt{MainActivity} que herda de \texttt{ComponentActivity}. Responda:

\begin{enumerate}[label=\alph*)]
    \item Explique qual dos dois trechos é utilizado durante o \textbf{desenvolvimento} (para pré-visualização na IDE) e qual é invocado quando o aplicativo é executado no \textbf{dispositivo real}.
    \item O que aconteceria se a função \texttt{LivrariaPreview()} fosse removida do código? O aplicativo pararia de funcionar no celular? Justifique.
\end{enumerate}

\vspace{0.3cm}
\resposta{6cm}

\vspace{0.5cm}

% ================================================================
% QUESTÃO 5 — Sensores
% ================================================================
\textbf{Questão 5 —} Dispositivos móveis modernos possuem diversos sensores embutidos, como Acelerômetro, Giroscópio, Magnetômetro, GPS, sensor de proximidade e sensor de luminosidade. Responda:

\begin{enumerate}[label=\alph*)]
    \item Explique a função de cada um dos três sensores: \textbf{Acelerômetro}, \textbf{Giroscópio} e \textbf{Magnetômetro}.
    \item Uma equipe pretende desenvolver um aplicativo de \textbf{nível digital} (bolha) para smartphones, que indica se uma superfície está perfeitamente nivelada. Qual sensor seria o principal para essa funcionalidade? Justifique.
    \item Em um aplicativo de \textbf{bússola digital}, qual sensor seria o mais adequado? Justifique.
\end{enumerate}

\vspace{0.3cm}
\resposta{7cm}

\vspace{0.5cm}

% ================================================================
% QUESTÃO 6 — Lambdas
% ================================================================
\textbf{Questão 6 —} Analise o seguinte trecho de código Kotlin:

\begin{lstlisting}[basicstyle=\small\ttfamily, breaklines=true]
val multiplicar = { x: Int, y: Int -> x * y }
val resultado = multiplicar(4, 7)
println(resultado)
\end{lstlisting}

\begin{enumerate}[label=\alph*)]
    \item Explique o que é uma \textbf{função lambda} (função anônima) em Kotlin e por que dizemos que funções são ``cidadãs de primeira classe'' nessa linguagem.
    \item No código acima, o que a variável \texttt{multiplicar} armazena e qual será o valor impresso pelo \texttt{println}?
    \item Reescreva o trecho acima criando uma lambda que receba uma \texttt{String} e retorne essa string em letras maiúsculas.
\end{enumerate}

\vspace{0.3cm}
\resposta{7cm}

\vspace{0.5cm}

% ================================================================
% QUESTÃO 7 — Coleções
% ================================================================
\textbf{Questão 7 —} Considere o trecho de código Kotlin abaixo:

\begin{lstlisting}[basicstyle=\small\ttfamily, breaklines=true]
val frutas = mutableListOf("Manga", "Uva", "Manga", "Abacaxi")
val conjunto = mutableSetOf("Manga", "Uva", "Manga", "Abacaxi")
val mapa = mutableMapOf("A" to 1, "B" to 2, "A" to 3)
\end{lstlisting}

\begin{enumerate}[label=\alph*)]
    \item Descreva as principais diferenças entre \texttt{List}, \texttt{Set} e \texttt{Map} em Kotlin.
    \item Qual será o valor de \texttt{frutas.size}, \texttt{conjunto.size} e \texttt{mapa.size}? Justifique cada resposta.
    \item Explique a diferença entre as variantes \texttt{mutableListOf} e \texttt{listOf}. A declaração com \texttt{val} impede a adição de novos elementos a uma coleção mutável?
\end{enumerate}

\vspace{0.3cm}
\resposta{7cm}

\vspace{0.5cm}

% ================================================================
% QUESTÃO 8 — filter, map, higher-order
% ================================================================
\textbf{Questão 8 —} Sobre as funções de alta ordem \texttt{filter} e \texttt{map} em Kotlin, responda:

\begin{enumerate}[label=\alph*)]
    \item Explique o que cada uma dessas funções faz e se elas alteram a coleção original ou retornam uma nova coleção.
    \item Qual é o resultado das expressões abaixo? Justifique.
\end{enumerate}

\begin{lstlisting}[basicstyle=\small\ttfamily, breaklines=true]
val numeros = listOf(1, 2, 3, 4, 5, 6)
val pares = numeros.filter { it % 2 == 0 }
val dobrados = numeros.map { it * 2 }
println(pares)
println(dobrados)
\end{lstlisting}

\begin{enumerate}[label=\alph*), resume]
    \item O que é o identificador implícito \texttt{it} utilizado dentro das lambdas acima?
    \item Escreva uma expressão que, a partir de \texttt{numeros}, retorne uma lista contendo apenas os valores maiores que 3, multiplicados por 10 (encadeando \texttt{filter} e \texttt{map}).
\end{enumerate}

\vspace{0.3cm}
\resposta{7cm}

\vspace{0.5cm}

% ================================================================
% QUESTÕES EXTRAS
% ================================================================
\hrule
\vspace{0.3cm}
\begin{center}
    \textbf{QUESTÕES EXTRAS DE REVISÃO}
\end{center}
\vspace{0.3cm}

% ================================================================
% QUESTÃO 9 — data class
% ================================================================
\textbf{Questão 9 —} No código de apoio, a classe \texttt{Livro} é declarada como uma \texttt{data class}. Explique o que diferencia uma \texttt{data class} de uma classe comum em Kotlin. Quais métodos são gerados automaticamente pelo compilador ao utilizar essa palavra-chave?

\vspace{0.3cm}
\resposta{5cm}

\vspace{0.5cm}

% ================================================================
% QUESTÃO 10 — Modifier
% ================================================================
\textbf{Questão 10 —} No Jetpack Compose, o conceito de \texttt{Modifier} é amplamente utilizado, como pode ser visto no código de apoio (ex: \texttt{Modifier.fillMaxWidth()}, \texttt{Modifier.padding(8.dp)}, \texttt{Modifier.height(180.dp)}). Responda:

\begin{enumerate}[label=\alph*)]
    \item O que é um \texttt{Modifier} no Jetpack Compose e qual a sua finalidade?
    \item Por que diversos \texttt{Modifiers} podem ser encadeados (ex: \texttt{Modifier.fillMaxWidth().padding(8.dp)})? Que padrão de projeto isso lembra?
\end{enumerate}

\vspace{0.3cm}
\resposta{6cm}

\vspace{0.5cm}

% ================================================================
% QUESTÃO 11 — LazyColumn vs Column
% ================================================================
\textbf{Questão 11 —} No código de apoio, a função \texttt{ListaLivros()} utiliza um \texttt{LazyColumn} para exibir os cards de livros. Explique a diferença entre \texttt{LazyColumn} e \texttt{Column} e em quais cenários o uso de \texttt{LazyColumn} é preferível.

\vspace{0.3cm}
\resposta{5cm}

\vspace{0.5cm}

% ================================================================
% QUESTÃO 12 — Exercício prático com coleções
% ================================================================
\textbf{Questão 12 —} Considerando a lista de livros abaixo em Kotlin:

\begin{lstlisting}[basicstyle=\small\ttfamily, breaklines=true]
data class Livro(
    val titulo: String,
    val autor: String,
    val preco: Double,
    val avaliacao: Int
)

val livros = listOf(
    Livro("Kotlin em Acao", "Dmitry", 89.90, 5),
    Livro("Clean Code", "Robert", 69.90, 4),
    Livro("O Hobbit", "Tolkien", 34.90, 5),
    Livro("Design Patterns", "GoF", 120.00, 3)
)
\end{lstlisting}

Escreva trechos de código Kotlin que realizem as seguintes operações utilizando funções de alta ordem:

\begin{enumerate}[label=\alph*)]
    \item Filtrar somente os livros com preço menor que R\$ 70,00.
    \item Obter uma lista contendo apenas os títulos dos livros com avaliação igual a 5.
    \item Calcular o preço total de todos os livros da lista (dica: use \texttt{sumOf}).
\end{enumerate}

\vspace{0.3cm}
\resposta{7cm}

\vspace{1cm}
\hrule
\vspace{0.3cm}
\begin{center}
    \textit{Boa revisão!}
\end{center}

\end{document}
